\documentclass[12pt]{article}
\usepackage[utf8]{inputenc}
\usepackage{amsmath}
\usepackage{amssymb}
\usepackage{graphicx}
\usepackage{booktabs}
\usepackage{listings}
\usepackage{xcolor}
\usepackage{hyperref}
\usepackage{geometry}
\geometry{margin=1in}

\lstset{
    language=Python,
    basicstyle=\ttfamily\small,
    keywordstyle=\color{blue}\bfseries,
    commentstyle=\color{green!60!black},
    stringstyle=\color{red},
    numbers=left,
    numberstyle=\tiny\color{gray},
    stepnumber=1,
    numbersep=5pt,
    backgroundcolor=\color{gray!10},
    frame=single,
    breaklines=true,
    breakatwhitespace=false,
    tabsize=4,
    showstringspaces=false
}

\title{Build Costs Calculation Documentation}
\author{}
\date{}

\begin{document}

\maketitle

\section{Introduction}

The \texttt{build\_costs.py} script calculates the construction costs for a transmission line project. It computes costs for conductors, structures, and converters, applies terrain adjustments via weighted miles, incorporates contingency factors, and calculates both regulatory (AFUDC) and societal (present value) perspectives.

\section{Method Overview}

\subsection{Purpose}

This script calculates the total build cost of a transmission line project by:
\begin{enumerate}
    \item Loading project technical details and constructing category via \texttt{calculation\_utils.build\_category\_string}
    \item Loading base costs from \texttt{10\_project\_category\_build\_costs.yaml} for that category
    \item Adjusting costs for terrain complexity using \texttt{weighted\_miles.calculate\_weighted\_miles()}
    \item Applying per-component contingency factors (conductor, structure, converter)
    \item Loading AFUDC setup via \texttt{financial\_utils.load\_afudc\_setup}; if enabled, calculating AFUDC capitalized cost using timing patterns for regulatory perspective
    \item Computing present value of build costs (spread evenly over construction period) for societal perspective
    \item Writing results via \texttt{CTCCOutputManager.add\_build\_costs} and \texttt{write\_batch\_summary}
\end{enumerate}

\subsection{Calculation Approach}

The build cost calculation follows this structure:

\begin{align}
\text{Conductor Cost} &= (\text{Variable Cost/Mile} \times \text{Weighted Miles}) + \text{Fixed Cost} \\
\text{Structure Cost} &= \text{Variable Cost/Mile} \times \text{Weighted Miles} \quad \text{(if not reconductoring)} \\
\text{Converter Cost} &= \text{Fixed Cost} \times \text{Number of Converters} \quad \text{(if not reconductoring)} \\
\text{Total Base Cost} &= \text{Conductor} + \text{Structure} + \text{Converter} \\
\text{Conductor with Conting.} &= \text{Conductor} \times (1 + \text{conductor\_contingency}) \\
\text{Structure with Conting.} &= \text{Structure} \times (1 + \text{structure\_contingency}) \\
\text{Converter with Conting.} &= \text{Converter} \times (1 + \text{converter\_contingency}) \\
\text{Total with Contingencies} &= \text{sum of above three}
\end{align}

For reconductoring projects, structure and converter costs are set to zero since existing infrastructure is reused.

\subsection{Inputs and Outputs}

\textbf{Inputs:}
\begin{itemize}
    \item Project technical details (construction type, AC/DC, capacity, conductor type, etc.)
    \item Physical details (terrain distribution)
    \item Build cost parameters from \texttt{10\_project\_category\_build\_costs.yaml} (by category from \texttt{build\_category\_string})
    \item Per-component contingency percentages (\texttt{conductor\_contingency}, \texttt{structure\_contingency}, \texttt{converter\_contingency}) from financing
    \item Financing parameters (WACC, inflation); AFUDC from \texttt{load\_afudc\_setup} (timing patterns, AFUDC rate)
\end{itemize}

\textbf{Outputs:}
\begin{itemize}
    \item Nominal costs (with contingencies)
    \item AFUDC capitalized costs (regulatory perspective)
    \item Present value costs (societal perspective)
    \item Component-level costs (conductor, structure, converter)
\end{itemize}

\subsection{Integration with Other Scripts}

This script integrates with:
\begin{itemize}
    \item \texttt{smart\_loaders.py} - \texttt{load\_project\_technical\_details}, \texttt{load\_physical\_details}, \texttt{load\_contingencies}, \texttt{load\_financing\_details}, \texttt{load\_cost\_timing\_patterns}, \texttt{load\_afudc\_config}, \texttt{get\_project\_data\_raw}, \texttt{get\_financing\_data\_raw}
    \item \texttt{weighted\_miles.py} - \texttt{calculate\_weighted\_miles} for terrain-adjusted miles
    \item \texttt{calculation\_utils.py} - \texttt{build\_category\_string} for project category
    \item \texttt{financial\_utils.py} - \texttt{calculate\_present\_value}, \texttt{calculate\_afudc\_capitalized\_cost}, \texttt{load\_afudc\_setup}, \texttt{calculate\_construction\_start\_year}, \texttt{validate\_discount\_rate}
    \item \texttt{path\_config.py} - \texttt{YAMLS\_DIR} for build costs YAML
    \item \texttt{smart\_output.py} - \texttt{CTCCOutputManager} (writes CSV or JSON via \texttt{add\_build\_costs}, \texttt{write\_batch\_summary})
\end{itemize}

\section{Key Functions}

\subsection{\texttt{load\_costs(category, total\_miles, number\_of\_converters, contingencies, reconductoring)}}

Loads and calculates build costs for the specified project category.

\textbf{Parameters:}
\begin{itemize}
    \item \texttt{category} (str): Project category identifier in format \texttt{"ConstructionType/AC\_DC/CapacityMW/ConductorType/ConverterType"}
    \item \texttt{total\_miles} (float): Total miles of transmission line
    \item \texttt{number\_of\_converters} (int): Number of converters needed (0 for AC)
    \item \texttt{contingencies} (dict): Dictionary with keys \texttt{"conductor\_contingency"}, \texttt{"structure\_contingency"}, \texttt{"converter\_contingency"}
    \item \texttt{reconductoring} (bool): Whether this is a reconductoring project
\end{itemize}

\textbf{Returns:} A \texttt{BuildCosts} dataclass with fields: \texttt{total\_cost}, \texttt{total\_cost\_with\_contingencies}, \texttt{conductor\_cost}, \texttt{structure\_cost}, \texttt{converter\_cost}, \texttt{conductor\_cost\_with\_contingencies}, \texttt{structure\_cost\_with\_contingencies}, \texttt{converter\_cost\_with\_contingencies}, \texttt{weighted\_miles}, \texttt{average\_terrain\_multiplier}. Structure and converter costs are 0 for reconductoring.

\textbf{Key Logic:}
\begin{lstlisting}
weighted_miles, average_terrain_multiplier = calculate_weighted_miles()
conductor_cost = (variable_conductor_cost_per_mile * weighted_miles) + fixed_conductor_cost
if reconductoring:
    structure_cost = 0
    converter_cost = 0
else:
    structure_cost = variable_structure_cost_per_mile * weighted_miles
    converter_cost = fixed_converter_cost * number_of_converters
total_cost = conductor_cost + structure_cost + converter_cost
conductor_cost_with_contingencies = conductor_cost * (1 + contingencies["conductor_contingency"])
structure_cost_with_contingencies = structure_cost * (1 + contingencies["structure_contingency"])
converter_cost_with_contingencies = converter_cost * (1 + contingencies["converter_contingency"])
total_cost_with_contingencies = (conductor_cost_with_contingencies
    + structure_cost_with_contingencies + converter_cost_with_contingencies)
return BuildCosts(...)
\end{lstlisting}

\subsection{\texttt{main()}}

Main execution function that orchestrates the build cost calculation and output.

\textbf{Workflow:}
\begin{enumerate}
    \item Loads project technical details; constructs category via \texttt{build\_category\_string(project\_details)}
    \item Determines \texttt{number\_of\_converters} from \texttt{get\_project\_data\_raw()["project"]["number\_of\_converters"]} if DC, else 0
    \item Loads physical details, contingencies, financing; calls \texttt{load\_costs(category, total\_miles, number\_of\_converters, contingencies, project\_details.reconductoring)}
    \item Loads AFUDC setup via \texttt{load\_afudc\_setup()}; if \texttt{apply\_afudc}, calculates AFUDC capitalized cost via \texttt{calculate\_afudc\_capitalized\_cost(...)}
    \item Computes \texttt{construction\_start\_year} = \texttt{calculate\_construction\_start\_year(delay\_years)}; if \texttt{construction\_years > 0}, \texttt{annual\_build\_cost} = total with contingencies / construction\_years, \texttt{build\_cost\_pv} = \texttt{calculate\_present\_value(annual\_build\_cost, wacc\_real, construction\_years, construction\_start\_year)}; else one-time PV at \texttt{construction\_start\_year}
    \item Builds \texttt{results} dict (\texttt{total\_nominal}, \texttt{total\_afudc}, \texttt{total\_pv}, component nominals and zeros for component AFUDC); \texttt{CTCCOutputManager().add\_build\_costs(results).write\_batch\_summary()}
\end{enumerate}

\section{Example}

The following example demonstrates a typical usage of the build costs calculation:

\begin{lstlisting}
# The script is typically called from ctcc.py, but can be run standalone
# It automatically loads configuration from YAML files

# Example: 500 MW Overhead AC line with Advanced Aluminum Conductor
# Category: "Overhead/AC/500MW/Advanced Aluminum/NA"
# 
# Inputs (from YAML):
#   - Construction type: Overhead
#   - AC/DC: AC
#   - Capacity: 500 MW
#   - Conductor: Advanced Aluminum
#   - Total miles: 100 miles
#   - Terrain: 60% flat, 40% forested
#   - Contingencies: 15% conductor, 20% structure, 10% converter
#
# Calculation:
#   1. Weighted miles = 100 * (0.6*1.0 + 0.4*1.2) = 108 miles
#   2. Conductor cost = ($500k/mile * 108) + $1M = $55M
#   3. Structure cost = $200k/mile * 108 = $21.6M
#   4. Converter cost = 0 (AC project)
#   5. Total base = $76.6M
#   6. With contingencies = $76.6M * (1 + weighted avg) ≈ $88M
#
# Output:
#   - Nominal cost: $88M
#   - AFUDC capitalized: $92M (if AFUDC enabled)
#   - Present value: $88M (discounted at real WACC)
\end{lstlisting}

\textbf{Integration Example:}

\begin{lstlisting}
# When called from ctcc.py:
# 1. Script loads project details from YAML
# 2. Calculates build costs
# 3. Writes to CSV via CTCCOutputManager:
csv_manager = CTCCOutputManager()
results = {
    "total_nominal": costs.total_cost_with_contingencies,
    "total_afudc": capitalized_cost_with_contingencies if afudc_setup.apply_afudc else 0,
    "total_pv": build_cost_pv,
    "conductor_nominal": costs.conductor_cost_with_contingencies,
    "structure_nominal": costs.structure_cost_with_contingencies,
    "converter_nominal": costs.converter_cost_with_contingencies,
    "conductor_afudc": 0, "structure_afudc": 0, "converter_afudc": 0,
}
csv_manager.add_build_costs(results)
csv_manager.write_batch_summary()
\end{lstlisting}

\section{Key Implementation Details}

\subsection{Reconductoring Handling}

For reconductoring projects, the script sets structure and converter costs to zero since existing infrastructure is reused. Only conductor costs are calculated.

\subsection{AFUDC Calculation}

If AFUDC is enabled (\texttt{load\_afudc\_setup().apply\_afudc}), the script calculates capitalized costs via \texttt{calculate\_afudc\_capitalized\_cost} using the AFUDC rate and cost timing patterns (\texttt{afudc\_setup.timing\_patterns["build\_costs"]}). This represents the regulatory perspective where construction financing costs are capitalized.

\subsection{Present Value of Build Costs}

The script computes present value of build costs spread evenly over the construction period: \texttt{annual\_build\_cost} = \texttt{total\_cost\_with\_contingencies} / \texttt{construction\_years}, then \texttt{calculate\_present\_value(annual\_build\_cost, wacc\_real, construction\_years, construction\_start\_year)}. If \texttt{construction\_years} is 0, build cost is treated as a one-time cost at \texttt{construction\_start\_year}.

\subsection{Weighted Miles}

The script uses \texttt{calculate\_weighted\_miles()} to adjust costs for terrain complexity. Different terrain types have different cost multipliers, and the weighted average is applied to the total miles.

\end{document}

